\documentclass{beamer}

\usepackage[utf8]{inputenc}
\usepackage[english]{babel}
\usepackage{textpos}
\usepackage{graphicx}
\usepackage{textcomp}
\usepackage{animate}
\usepackage{comment}
%\usepackage{multicol}
%\setlength{\columnsep}{5cm}
\usepackage{amsmath}
\usepackage{amsfonts}
\usepackage{verbatim}
\usepackage{varwidth}
\usepackage{caption}
\usepackage{dcolumn}
\usepackage{amssymb}
\usepackage{algorithm}
\usepackage{algorithmic}
%\usepackage{physics}
\usepackage{hyperref}
\usepackage{scalerel,stackengine}
%\usepackage{fancyvrb}
%\usepackage{enumitem}


\newcommand{\code}[1]{\texttt{#1}}
\newcommand{\shp}{$>>>$ }
\newcommand{\tab}{\hspace*{22pt}}

\usetheme{Madrid}

\title{Chapter 10: \\ Classes, Objects, and Methods: \\ Object-Oriented Programming w/ Python}
\author{EEC 3150}
\date{}

\institute{Trevecca Nazarene University}
 


%%%%%%%%% BEGIN DOCUMENT %%%%%%%%%
%%%%%%%%% BEGIN DOCUMENT %%%%%%%%%
%%%%%%%%% BEGIN DOCUMENT %%%%%%%%%

\begin{document}

% get title page 
\frame{\titlepage}
 
%%%%%%%%%%
%%%%%%%%%%
\begin{frame}	
	\frametitle{Outline}
	\tableofcontents	
\end{frame}


%%%%%%%%%%
%%%%%%%%%%
\section{Classes, Objects, and Methods}

\begin{frame}	
	\center{\Huge{Classes, Objects, and Methods}}
\end{frame}

\begin{frame}	
	\frametitle{Classes, Objects, and Methods}
	
	\begin{block}{Data types vs. objects:}
		\minipage{0.35\textwidth}
			Definitions:
			\begin{itemize}
				\item A \textbf{data type} is a particular format for representing data of a certain kind
				\item An \textbf{object} is an particular instance of a data type
			\end{itemize}
		\endminipage\hfill
		\minipage{0.65\textwidth}
			Examples:
			\begin{itemize}
				\item data type: \code{int} \\
						object: \code{1}
				\item data type: \code{str} \\
						object: \code{``hello''}
				\item data type: \code{builtin\_function\_or\_method} \\
						object: \code{print()}
			\end{itemize}
		\endminipage\hfill
	\end{block}
	
\end{frame}

\begin{frame}	
	\frametitle{Classes, Objects, and Methods}
	
	\begin{alertblock}{Defn: classes and objects}
		A \textbf{class} is (a template for creating) an abstract data type
		\begin{itemize}
			\item We can create our own custom data types
			\item All Python data types are actually classes
		\end{itemize}
		An \textbf{object} is an particular instance of a class
		
		\vspace{10pt}
		Analogy:
		\begin{itemize}
			\item class: car \\
					object: Honda Civic
		\end{itemize}
	\end{alertblock}
	
\end{frame}

\begin{frame}	
	\frametitle{Classes, Objects, and Methods}
	
	\begin{block}{First class: \code{Person}}
		\begin{scriptsize}
		Define a \code{Person} class with a \code{name} attribute:
		
		\vspace{15pt}
		\minipage{0.38\textwidth}
			Defining the \code{Person} class:
		
			\vspace{10pt}
			\code{class Person: \\
					\tab def \_\_init\_\_(self, name):
					\tab \tab self.\_name = name
			}
			
			\rule[0pt]{\linewidth}{1pt}
		
			\vspace{10pt}
			Creating a \code{Person} object:
			
			\vspace{10pt}
			\code{p = Person(``Graham'')
			}
			
			\rule[0pt]{\linewidth}{1pt}
		
			\vspace{10pt}
			Getting \code{p}'s name:
			
			\vspace{10pt}
			\code{p.\_name
			}
		\endminipage\hfill
		\minipage{0.62\textwidth}
			\begin{itemize}
%				\setlength\itemsep{-2pt}
				\item Create a class by using the \code{class} keyword, the name of the class, and a colon
				\item \emph{method}: any function defined in the scope of a class
				\begin{itemize}
					\begin{scriptsize}
%					\setlength\itemsep{0pt}
					\item Every class should have a \emph{constructor} method called \code{\_\_init\_\_} which takes \code{self} and any attributes as arguments
					\item It is automatically called whenever a new class is created
					\item The \code{self} variable is a reference to whatever particular object is calling a method; all methods definitions must have it, but nothing is ever passed to it when called
					\end{scriptsize}
				\end{itemize}
				\item \emph{attribute}: any variable associated with a class
				\begin{itemize}
					\begin{scriptsize}
%					\setlength\itemsep{0pt}
					\item Set any attributes passed to the constructor
					\item It is convention for all attributes to start with underscores
					\end{scriptsize}
				\end{itemize}
			\end{itemize}
		\endminipage\hfill
		\end{scriptsize}
	\end{block}
	
\end{frame}

\begin{frame}	
	\frametitle{Classes, Objects, and Methods}
	
	\begin{block}{First class: \code{Person}}
		\begin{scriptsize}
		Define a \code{Person} class with a \code{name} attribute:
		
		\vspace{15pt}
		\minipage{0.38\textwidth}
			Defining the \code{Person} class:
		
			\vspace{10pt}
			\code{class Person: \\
					\tab def \_\_init\_\_(self, name):
					\tab \tab self.\_name = name
			}
			
			\rule[0pt]{\linewidth}{1pt}
		
			\vspace{10pt}
			Creating a \code{Person} object:
			
			\vspace{10pt}
			\code{p = Person(``Graham'')
			}
			
			\rule[0pt]{\linewidth}{1pt}
		
			\vspace{10pt}
			Getting \code{p}'s name:
			
			\vspace{10pt}
			\code{p.\_name
			}
		\endminipage\hfill
		\minipage{0.62\textwidth}
			\begin{itemize}
%				\setlength\itemsep{-2pt}
				\item To create a \code{Person} object, simply pass all necessary attributes to the \code{Person} function
				\begin{itemize}
					\begin{scriptsize}
%					\setlength\itemsep{0pt}
					\item \code{Person} is the class, \code{p} is the object
					\item Again, nothing is passed for the \code{self} argument
					\end{scriptsize}
				\end{itemize}
				\item To access or change an object's attributea or call a method on it, use dot notation
			\end{itemize}
		\endminipage\hfill
		\end{scriptsize}
	\end{block}
	
\end{frame}

\begin{frame}	
	\frametitle{Classes, Objects, and Methods}
	
	\begin{block}{\code{Person} with multiple attributes and methods:}
		\begin{scriptsize}
		\minipage{0.45\textwidth}
			\code{class Person: \\
					\tab def \_\_init\_\_(self, name, age): \\
					\tab \tab self.\_name = name \\
					\tab \tab self.\_age = age \\
					\tab def HaveBirthday(self): \\
					\tab \tab self.\_age += 1 \\
					\tab \\
					p = Person(``Bill'', 50) \\
					p.HaveBirthday()
			}
		\endminipage\hfill
		\minipage{0.55\textwidth}
			\begin{itemize}
%				\setlength\itemsep{-2pt}
				\item \code{Person} class has two attributes (\code{\_name,\_age}) and two methods (\code{\_\_init\_\_,HaveBirthday})
				\item \code{HaveBirthday} will increment \code{\_age} by 1
			\end{itemize}
		\endminipage\hfill
		\end{scriptsize}
	\end{block}
	
	\begin{block}{The \code{vars()} and \code{isinstance()} functions:}
		\begin{itemize}
			\item \code{vars(x)} returns a dictionary of all attributes of the object \code{x}
			\item \code{isinstance(x,Class)} check whether the object \code{x} is an instance of the \code{Class} class
		\end{itemize}
	\end{block}
	
\end{frame}

\begin{frame}	
	\frametitle{Classes, Objects, and Methods}
	
	\begin{exampleblock}{\code{Rectangle} class:}
		Create a \code{Rectangle} class with the following attributes
		\code{
		\begin{itemize}
			\item length
			\item width
		\end{itemize}
		}
		and the following methods
		\code{
		\begin{itemize}
			\item get\_area()
			\item get\_perimeter()
		\end{itemize}
		}
	\end{exampleblock}
	
\end{frame}





\end{document}




